\section{One selected nonlinear dynamics for both source constructions}
\label{cp:selected-section}

The sparse-metric lemma first supplies exact geometry on any finite source
space. The transfer proposition then constructs and analyzes one autonomous
selected model, with its optimizer and comparison coefficients local to that
proof. Two later propositions supply its source spaces from whole-sphere
polynomials or finite-panel Taylor expansions. They use only the transfer
statement's error and storage interface.

\begin{lemma}[Sparse coordinates with an exact source metric]
\label{cp:selection}
Let $E\subset\R^n$ be a linear subspace containing $\mathbf1_n$,
with $r=\dim E$.
There are a set $I$ of at most $9r$ coordinates, a positive diagonal
matrix $D$ on $I$, and a positive definite matrix $M$ such that
\[
 u_I^\top Mv_I=u^\top v/n\quad(u,v\in E),\qquad
 D/4\preceq M\preceq D,\qquad
 \mathbf1^\top M\mathbf1=1,\quad\mathbf1^\top D\mathbf1\le4.
\]
Write $\|u\|_M=(u^\top Mu)^{1/2}$. Consequently a coordinate
multiplication operator $B$ satisfies
$\|B\|_{M\to M}\le2\max_i|B_{ii}|$, and an activation with
$|\phi'|\le s$, $|\phi(0)|\le b$ satisfies
$\|\phi(u)-\phi(v)\|_M\le2s\|u-v\|_M$ and
$\|\phi(u)\|_M\le2b+2s\|u\|_M$.
\end{lemma}
\begin{proof}
We give the finite sparsification argument. For vectors $v_i\in\R^r$
with $\sum_i v_iv_i^\top=I$, start with $A=0$, lower barrier $l=-3r$
and upper barrier $u=6r$. Maintain
\[
 lI\prec A\prec uI,\qquad
 \Phi_l(A)=\operatorname{tr}(A-lI)^{-1}\le\tfrac13,
 \quad\Phi^u(A)=\operatorname{tr}(uI-A)^{-1}\le\tfrac16.
\]
Set $l'=l+1$, $u'=u+2$, and define
\[
 U(v)=\frac{v^\top(u'I-A)^{-2}v}{\Phi^u(A)-\Phi^{u'}(A)}
       +v^\top(u'I-A)^{-1}v,\qquad
 L(v)=\frac{v^\top(A-l'I)^{-2}v}{\Phi_{l'}(A)-\Phi_l(A)}
       -v^\top(A-l'I)^{-1}v.
\]
The lower potential implies that every eigenvalue is at least three
above $l$, so both inverse matrices and both denominators are positive.
The identity
$(B+t vv^\top)^{-1}=B^{-1}-B^{-1}vv^\top B^{-1}/(t^{-1}+v^\top B^{-1}v)$
shows that $U(v)\le t^{-1}\le L(v)$ preserves the two potential
inequalities at the shifted barriers. For the upper potential, the
positive first term of $U(v)$ also preserves the strict spectral barrier.

Such a nonzero $v_i$ exists. If $a_j$ are the eigenvalues of $A$, then
$(u-a_j)^{-1}-(u'-a_j)^{-1}\ge2(u'-a_j)^{-2}$, hence
$\sum_iU(v_i)\le1/2+\Phi^{u'}(A)\le2/3$.
For the lower sum put $a'_j=(a_j-l')^{-1}$,
$b'_j=(a_j-l)^{-1}=a'_j/(1+a'_j)$ and
$\Delta=\sum_j a'_jb'_j=\Phi_{l'}-\Phi_l$.
Weighted Cauchy--Schwarz gives
\[
 \Delta^2\le\Big(\sum_j(a'_j)^2b'_j\Big)\sum_jb'_j
 \le\sum_j(a'_j)^2-\Delta.
\]
Therefore $\sum_iL(v_i)\ge1-\Phi_l(A)\ge2/3$.
Choose $t>0$ with reciprocal in the resulting nonempty interval and
add $t v_iv_i^\top$. After $9r$ iterations the barriers are $6r,24r$.
Dividing all accumulated weights by $6r$ gives nonnegative weights
$s_i$, supported on at most $9r$ indices, with
$I\preceq\sum_i s_i v_iv_i^\top\preceq4I$.

Choose $U$ with range $E$ and $U^\top U/n=I$, apply this result to
$v_i=U_{i,:}^\top/\sqrt n$, and set
$P=U_I$, $D=\operatorname{diag}(s_i/n)$, $G=P^\top DP$ and
$Z=D^{1/2}P$. Define
\begin{equation}
 M=D^{1/2}\{ZG^{-2}Z^\top+I-ZG^{-1}Z^\top\}D^{1/2}.
 \label{cp:metric}
\end{equation}
On $\operatorname{ran}Z$ the middle matrix has eigenvalues those of
$G^{-1}$, and on its orthogonal complement it is the identity.
Thus $D/4\preceq M\preceq D$, while multiplication gives $P^\top MP=I$.
The constant source proves the two mass assertions. Finally
$\|Bu\|_M\le\|Bu\|_D\le\max_i|B_{ii}|\|u\|_D\le
2\max_i|B_{ii}|\|u\|_M$. Apply this to coordinate divided differences
and use $\|\mathbf1\|_M=1$ for the activation claims.
\end{proof}

\subsection{A portable source-to-runtime construction}

\begin{proposition}[Source-to-runtime transfer]
\label{cp:selected}
Work on the fitting and analytic-source event of \Cref{cp:fit,cp:source}.
Write $W_0^{(j)}=W^{(j)}(0)$ for the initialized matrices in this statement.
Fix
\[
 T=32\frac m\gamma\log(en),\qquad
 0<\eta\le\min\{1,Y,16Ym/\gamma\},\qquad
 q\ge\max\{m,d\}.
\]
For each layer let $E_j\subset\R^n$ have dimension at most $q$ and contain
the constant, exact initialized training features and their initialized
forward images when present, and (at the first layer) the columns of
$W_0^{(1)}$. Suppose these spaces approximate the families
\eqref{cp:source-families} to coordinate accuracy $\eta$ on $[0,T]$,
with forward queries in $\mathcal X$ and backward queries at training inputs.
Require exact initialized-image pairing: each image approximant is the
corresponding initialized map applied to its base approximant.
There is an autonomous selected nonlinear predictor satisfying
\begin{align}
 \|f_{\rm model}-f_n\|_{\mathcal X}
 &\le C\frac{Ym}\gamma\left(1+\sqrt{\frac m\gamma}\right)
       (1+\sqrt{\log(en)})e^{\sqrt{\log(en)}}\eta
 \nonumber\\
 &\quad+C\frac{Ym}\gamma\left(1+\sqrt{\frac m\gamma}\right)
                      e^{-\gamma T/(4m)}.
 \label{cp:smallcap-comparison}
\end{align}
It fits, has a fitted limit, and retains at most
\begin{equation}
 1020(L+1)q^2+36q+10m(d+1)
 \label{cp:selected-inventory}
\end{equation}
real coordinates, including data, selected indices, metrics, initial copies
and working arrays. Here $C\ge1$.
In particular, the permissible choice
\begin{equation}
 \eta=\min\left\{1,Y,\frac{16Ym}\gamma,
 \frac{\gamma}{2Cmn(1+\sqrt{m/\gamma})
          (1+\sqrt{\log(en)})e^{\sqrt{\log(en)}}}\right\}
 \label{cp:selected-tolerance}
\end{equation}
gives error at most $Y/n$ eventually. Source bases, dense coefficients and
compilation scratch are discarded after the selected initial arrays are formed.
\end{proposition}
\begin{proof}
All auxiliary quantities in this proof belong to this construction.
Put $\lambda=\gamma/m$, $z=Y/\lambda$, $\alpha=Y/\sqrt\lambda$,
$S=16z$, and $\ell=\log(en)$, so $0<\lambda\le1$ and
$0<z\le\beta^{-30L}$. For activation estimates use the envelopes
$b=\beta-1$ and $s=t_2=\beta$; thus $|\phi_j(0)|\le b$,
$|\phi_j'|\le s$ and $|\phi_j''|\le t_2$ on the real line.
Write $c_n=(y_a-f_n(t,x_a))_{a=1}^m$ for the dense training deficits and
$\rho_n=\|c_n\|_2/\sqrt m$ for their RMS.
The public source bounds permit the local envelopes
\[
 u_0=\beta^{3L},\qquad \tau=\beta^{4L},\qquad
 K_{\rm src}=\beta^{21L}.
\]
Dense feature norms are at most $u_0$ and dense training response norms
at most $S\tau$ in neuron RMS. Dense training carriers on $[0,T]$
are the local fields
$k_{n,a}^{(L)}=w$ and
$k_{n,a}^{(j)}=W^{(j+1)\top}\delta_{n,a}^{(j+1)}$; by
\eqref{cp:carrier} they are bounded coordinatewise by
$32zK_{\rm src}\sqrt\ell$.
Dense vector norms below are normalized by $\sqrt n$; selected norms
are specified by their metrics. Every query $v$ denotes $x/\sqrt d$
for an input $x$ in the promised domain.

\paragraph{Autonomous selected model.}

Apply \Cref{cp:selection} to obtain $(I_j,M_j,D_j)$ and let
$q_j=|I_j|\le9q$. The selected moving variables are
$W_C^{(1)}\in\R^{q_1\times d}$,
$B_C^{(j)}\in\R^{q_j\times q_{j-1}}$ ($2\le j\le L$),
$w_C\in\R^{q_L}$ and $c_C\in\R^m$.
For $v=x/\sqrt d$ the forward pass is
\[
 z_C^{(1)}=W_C^{(1)}v,\qquad
 z_C^{(j)}=B_C^{(j)}h_C^{(j-1)},\qquad
 h_C^{(j)}=\phi_j(z_C^{(j)}).
\]
All layer norms and adjoints below use $M_j$; in particular
$(B_C^{(j)})^*=M_{j-1}^{-1}B_C^{(j)\top}M_j$.
Let $V_C$ have columns $h_C^{(L)}(v_a)/\sqrt m$ and define
\begin{equation}
 Q_C=V_C^*V_C,\qquad
 \widehat w_C=w_C+V_CQ_C^{-1}
       \{(y-c_C)/\sqrt m-V_C^*w_C\},\qquad
 f_C(v)=\langle\widehat w_C,h_C^{(L)}(v)\rangle_{M_L}.
 \label{cp:selected-readout}
\end{equation}
The statement's predictor is $f_{\rm model}(t,x)=f_C(t,x/\sqrt d)$;
physical time is suppressed in the local forward formulas.
Thus $f_C(v_a)=y_a-c_{C,a}$ exactly. Define training signals by
$k_{C,a}^{(L)}=\widehat w_C$,
$\delta_{C,a}^{(j)}=\phi_j'(z_{C,a}^{(j)})\odot k_{C,a}^{(j)}$ and
$k_{C,a}^{(j)}=(B_C^{(j+1)})^*\delta_{C,a}^{(j+1)}$.
The unnormalized sample matrix is
\begin{align}
 K_{C,ab}={}&\langle h_{C,a}^{(L)},h_{C,b}^{(L)}\rangle_{M_L}
 +\langle\delta_{C,a}^{(1)},\delta_{C,b}^{(1)}\rangle_{M_1}v_a^\top v_b
 \nonumber\\
 &+\sum_{j=2}^L\langle\delta_{C,a}^{(j)},\delta_{C,b}^{(j)}\rangle_{M_j}
                  \langle h_{C,a}^{(j-1)},h_{C,b}^{(j-1)}\rangle_{M_{j-1}}.
 \label{cp:selected-gram}
\end{align}
The autonomous optimizer is
\begin{align}
 \dot W_C^{(1)}&=\frac2m\sum_a c_{C,a}\delta_{C,a}^{(1)}v_a^\top,
 &\dot B_C^{(j)}&=\frac2m\sum_a c_{C,a}\delta_{C,a}^{(j)}
                                      h_{C,a}^{(j-1)\top}M_{j-1},
 \nonumber\\
 \dot w_C&=\frac2m\sum_a c_{C,a}h_{C,a}^{(L)},
 &\dot c_C&=-2K_Cc_C/m.
 \label{cp:selected-flow}
\end{align}
These are prescribed directions; coordinate gates need not be self-adjoint
in $M_j$, and no gradient interpretation of the corrected predictor is used.

For bases $U_j$ of $E_j$ with $U_j^\top U_j/n=I$, set $P_j=(U_j)_{I_j}$.
Initialize
\begin{equation}
 W_C^{(1)}(0)=(W_0^{(1)})_{I_1},\quad w_C(0)=0,\quad c_C(0)=y,
 \qquad
 B_C^{(j)}(0)=P_j\frac{U_j^\top W_0^{(j)}U_{j-1}}n
                         P_{j-1}^\top M_{j-1}.
 \label{cp:selected-initial}
\end{equation}
Since $P_j^\top M_jP_j=I$, these operators have norms at most eight.
They reproduce an initialized forward image exactly whenever both members
belong to the source spaces; their metric adjoints reproduce paired reverse
images. The exact additions therefore preserve the entire initialized
training forward pass and Gram. All bases, dense arrays and source
coefficients are discarded after these finite initial arrays are formed.

\begin{claim}[Independent fitting and endpoint tails]
With the notation of this construction, the following local bounds hold.
The selected dynamics exist for every $t\ge0$, have $Q_C\succeq\lambda I/4$,
fit the training labels and converge in parameter space. Write
\[
 H_1^c=2b+16s,\quad H_j^c=2b+18sH_{j-1}^c,\quad H_c=H_L^c,
 \quad d_j^c=2s(18s)^{L-j},\quad
 F_c=(d_1^c)^2+4H_c^2\sum_{j=2}^L(d_j^c)^2.
\]
For the selected parameter tuple
$\theta_C=(W_C^{(1)},(B_C^{(j)})_{j=2}^L,w_C)$ use
\[
 \|\theta_C\|_{\rm par}^2=
 \operatorname{tr}(W_C^{(1)\top}M_1W_C^{(1)})+
 \sum_{j=2}^L\|M_j^{1/2}B_C^{(j)}M_{j-1}^{-1/2}\|_F^2+
 \|w_C\|_{M_L}^2.
\]
Then $\|h_C^{(j)}\|\le2H_c$, $\|\widehat w_C\|\le5\alpha$,
$\rho_C:=\|c_C\|/\sqrt m\le Ye^{-\lambda t/2}$ and
$\int_0^\infty\|\dot\theta_C\|_{\rm par}\le2\alpha$.
There is a coefficient $\mathcal D\le C z(1+\lambda^{-1/2})$ such that
comparison on $[0,T]$ extends to all times and the fitted endpoint with
additional error at most $\mathcal D e^{-\lambda T/4}$.
Here and below $C$ depends only on activations and depth.
\end{claim}
\begin{proof}[Proof of the claim]
Stop at operator norm nine, feature norm $2H_c$, or Gram margin $\lambda/4$.
Expanding the velocities, with all cross-sample terms, gives
\begin{equation}
 -\frac d{dt}\rho_C^2=4c_C^\top K_Cc_C/m^2
       =\|\dot\theta_C\|_{\rm par}^2,
 \qquad K_C/m\succeq Q_C.
 \label{cp:selected-energy}
\end{equation}
Thus $-\dot\rho_C\ge\lambda\rho_C/2$ and
$\|\dot\theta_C\|\le2(-\dot\rho_C)/\sqrt\lambda$.
Integration gives the asserted residual and path bounds, as well as
$\int\rho_C\le2z$ and $\|w_C\|\le2\alpha$.
With $T_C=V_CQ_C^{-1}$ and $P_C=T_CV_C^*$, the readout is the
orthogonal sum $(I-P_C)w_C+T_C(y-c_C)/\sqrt m$.
Its squared norm is at most $4\alpha^2+16\alpha^2$.

Set $R_c=5\alpha$, $U_1^c=d_1^c$, $U_j^c=2H_cd_j^c$ for $j\ge2$,
and $F_1^c=2sU_1^c$, $F_j^c=2s(2H_cU_j^c+9F_{j-1}^c)$.
The latter recurrence has $F_L^c=F_c$. Backward propagation bounds
$\|\delta_C^{(j)}\|\le d_j^cR_c$, and integration bounds each hidden
parameter displacement by $20U_j^cY^2/\lambda^{3/2}$ and each feature
displacement by $20F_j^cY^2/\lambda^{3/2}$.
Under the stated cap, $H_c\le\beta^{3L}$, $F_c\le\beta^{15L}$ and
$z\le(16H_c\sqrt{F_c})^{-1}$; hence the last bound is at most
$5\sqrt\lambda/(64H_c^2)<\min\{1/(8H_c),\sqrt\lambda/8\}$.
The operator and feature stops improve strictly. The normalized training
feature matrix starts with smallest singular value at least
$\sqrt{\lambda/2}$ and moves by less than $\sqrt\lambda/8$, so the
Gram stop improves too. Smoothness on this strict Gram-positive domain
and bounded parameters give global existence. Finite path length and
\eqref{cp:selected-readout} give convergence; $c_C\to0$ gives fitting.

For quantitative tails, differentiating the inverse gives
\[
 \dot T_C=(I-P_C)\dot V_CQ_C^{-1}-T_C\dot V_C^*T_C,\qquad
 \dot P_C=(I-P_C)\dot V_CT_C^*+T_C\dot V_C^*(I-P_C).
\]
Thus $\|\dot T_C\|\le8\|\dot V_C\|/\lambda$,
$\|\dot P_C\|\le4\|\dot V_C\|/\sqrt\lambda$ and
$\|\dot V_C\|\le2R_cF_c\rho_C$.
Put $G_c=4H_c^2+R_c^2F_c$ and
\[
 B_w=4H_c+48R_cF_cz+4G_c/\sqrt\lambda,
 \qquad B_f=2H_cB_w+2R_c^2F_c.
\]
Differentiating the orthogonal readout formula proves
$\|\dot{\widehat w}_C\|\le B_w\rho_C$ and
$|\dot f_C(v)|\le B_f\rho_C$ uniformly over unit inputs.
Therefore $|f_C(\infty,v)-f_C(t,v)|\le2B_fz e^{-\lambda t/2}$.
For the dense tail, the RMS bounds hold at every real time, not merely
on the finite complex source domain. Indeed \Cref{cp:fit} bounds all
normalized operators by nine and gives
$\|w\|/\sqrt n\le2Y/\sqrt\lambda\le Su_0$.
The all-time RMS bounds of \Cref{cp:source}, or the same real
forward/backward induction, give
$\|h^{(j)}\|/\sqrt n\le u_0$ and
$\|\delta^{(j)}\|/\sqrt n\le S\tau$ for all real times.
The flow equations and sample Cauchy--Schwarz therefore imply
$|\dot f_n(v)|\le2\mathcal K\rho_n$, where
\[
 \mathcal K=u_0^2+S^2\tau^2[1+(L-1)u_0^2].
\]
Comparing each flow with its value at $T$ proves the claimed extension with
$\mathcal D=z(16\mathcal K+4B_f)$.
The small cap gives $\mathcal K\le C$, $B_f\le C(1+\lambda^{-1/2})$,
which establishes the stated parameter dependence.
\end{proof}

\paragraph{Comparison coefficients.}

Here is an explicit coefficient ledger for the comparison. Its size is
independent of the number of sources. Use the local envelopes
$u_0,\tau,K_{\rm src}$ fixed at this proof's start, and put
\[
 H_r=u_0+3,\quad H_C=2H_c,\quad
 P_h=6u_0+9,\quad P_\delta=6\tau+9,
\]
\[
 A_f=18+S^2(\tau+3)P_h,\qquad
 A_b=18+S^2H_rP_\delta.
\]
Define the forward-error recursion and the reference-response bounds by
\begin{align*}
 F_1^z&=1,\quad F_1^h=2s,\\
 F_j^z&=9F_{j-1}^h+H_r+A_f,\quad
 F_j^h=2sF_j^z\quad(j\ge2),\quad F=\max_jF_j^h,\\
 d_j^E&=2s(9s)^{L-j}+3H_c,\quad J_C=5\sqrt{F_c},\\
 J_R&=\left[(d_1^E)^2+H_r^2\sum_{j=2}^L(d_j^E)^2\right]^{1/2},\\
 D_s&=P_h+16zP_\delta\{1+(L-1)H_r^2\}
                   +(L-1)(16z)^2\tau^2P_h.
\end{align*}
The backward-error recursion is
\begin{align*}
 B_L&=2s(1+6zF+32zP_h)+2t_2F_L^z,\\
 B_j&=18sB_{j+1}+2s(d_{j+1}^EzH_C+A_b)+2t_2F_j^z
                       \quad(j<L),\\
 B_h&=\sqrt L\{H_C\max_jB_j+(\max_jd_j^E)zH_CF\},\\
 K_1&=12F+2B_h+20z(J_C+J_R)B_h+20D_s.
\end{align*}
Increasing any of these positive coefficients is harmless. Finally set
\begin{align}
 \mathcal B_n={}&400z^2F_c+2zK_1+40zF
                         +64z^2K_1K_{\rm src}\sqrt\ell,\nonumber\\
 \mathcal A_n={}&[H_C(1+6zF)+3\alpha F]
             (1+\lambda^{-1/2})(e^{\mathcal B_n}-1)\nonumber\\
 &+6H_CzF+32H_CzP_h/\sqrt\lambda+3\alpha F+16zP_h.
 \label{cp:selected-coefficients}
\end{align}

\begin{claim}[Finite-horizon comparison]
With the local coefficients \eqref{cp:selected-coefficients},
\[
 \sup_{0\le t\le T}\sup_{x\in\mathcal X}
 |f_C(t,x/\sqrt d)-f_n(t,x)|\le\mathcal A_n\eta.
\]
\end{claim}
\begin{proof}[Proof of the comparison claim]
We record both the source-energy estimate and the cancellation that
make the parameter dependence possible. For one layer abbreviate
$E=E_j$, $I=I_j$ and $M=M_j$. If $p\in E$ approximates $u$
to coordinate accuracy $e_u$, exact isometry and the mass bound give
\[
 \|u-p\|_n\le e_u,\quad \|u_I-p_I\|_M\le2e_u,
 \quad\|u_I\|_M\le\|u\|_n+3e_u.
\]
Expanding two inner products about $p,q\in E$ therefore gives
\begin{equation}
 |\langle u_I,v_I\rangle_M-\langle u,v\rangle_n|
 \le3e_u\|v\|_n+3e_v\|u\|_n+9e_ue_v.
 \label{cp:pair-defect}
\end{equation}
Thus feature-pair defects are at most $P_h\eta$, and response-pair
defects at most $SP_\delta\eta$; here $\eta\le1,S$ is used.

Let the proof reference have first layer and readout equal to the
coordinate restrictions of the dense ones, and hidden operators
\[
 B_R^{(j)}(t)=B_C^{(j)}(0)+\frac2m\int_0^t\sum_a
 c_{n,a}(u)(\delta_{n,a}^{(j)}(u))_{I_j}
              (h_{n,a}^{(j-1)}(u))_{I_{j-1}}^\top M_{j-1}\,du.
\]
Its preactivation, feature and response symbols below mean the restricted
dense values, not a fresh forward pass; its carrier symbols
mean $k_{R,a}^{(j)}=(k_{n,a}^{(j)})_{I_j}$ as well.
The exact initialized
pairing and operator bound eight make an initialized action defect at
most $18\eta$. In each learned action, subtract the dense outer-product
integral and use \eqref{cp:pair-defect} together with
$\int\rho_n\le2z$, $\|\delta_n^{(j)}\|_n\le S\tau$ and
$\|h_n^{(j)}\|_n\le u_0$. The resulting forward and reverse defects
are at most $A_f\eta$ and $A_b\eta$. The identical calculation for
the integrated raw readout gives
\begin{equation}
 |\langle w_R,h_R^{(L)}(v)\rangle-f_n(v)|\le16zP_h\eta.
 \label{cp:readout-observation}
\end{equation}

Write $V_n,V_R$ for the dense and restricted top feature maps, with
columns divided by $\sqrt m$. Their common approximants and exact
isometry give, for every $a\in\R^m$,
$\|V_Ra\|\le\|V_na\|+3\eta\|a\|$.
For $\nu=\|V_Rc_n/\sqrt m\|$, the dense energy identity gives
\begin{equation}
 \int_0^T\nu\le\alpha+6\eta z\le2\alpha.
 \label{cp:source-energy}
\end{equation}
Here let $\theta_n=(W^{(1)}/\sqrt n,W^{(2)},\ldots,W^{(L)},w/\sqrt n)$
with its Euclidean block norm. Indeed
$2\|V_nc_n/\sqrt m\|\le\|\dot\theta_n\|$,
whose integral is at most $2\alpha$.
Integrating source approximants for $\dot w_n=2V_nc_n/\sqrt m$
gives a vector in $E_L$ at coordinate error at most $4\eta z$.
Consequently $\|w_R\|\le2\alpha+12\eta z\le3\alpha$.
Likewise the restricted responses have norm at most $d_j^E\alpha$.
These arguments never differentiate an approximation error.

Use the hidden-parameter Hilbert norms of the preceding local fitting claim,
and let $\theta_{h,C},\theta_{h,R}$ denote the selected and reference tuples
of first-layer and hidden-layer matrices. Let
$\mathcal J_C,\mathcal J_R$ have sample columns, divided by $\sqrt m$,
equal respectively to the tuples
\[
 (\delta^{(1)}v_a^\top,
       (\delta^{(j)}h_a^{(j-1)\top}M_{j-1})_{j=2}^L).
\]
Thus $K_C/m=V_C^*V_C+\mathcal J_C^*\mathcal J_C$ and
$\dot\theta_{h,C}=2\mathcal J_Cc_C/\sqrt m$; the reference satisfies
the latter identity with $R,n$. Moreover
$\|\mathcal J_C\|\le J_C\alpha$,
$\|\mathcal J_R\|\le J_R\alpha$.
Put
\[
 a_h=\|\theta_{h,C}-\theta_{h,R}\|,\quad
 e=(c_C-c_n)/\sqrt m,\quad u=\|e\|,\quad z_w=w_C-w_R,
\]
\[
 T_C=V_CQ_C^{-1},\quad P_C=T_CV_C^*,\quad
 p=T_Ce,\quad \zeta=z_w+p,\quad b_e=\|p\|+\|\zeta\|,
 \quad E=a_h+b_e.
\]
All these error variables vanish at zero. Forward induction, first
subtracting the reference action and then the nonlinear gate, gives
\begin{equation}
 \|h_C^{(j)}-h_R^{(j)}\|\le F_j^h(a_h+\eta),
 \qquad \|V_C-V_R\|\le F(a_h+\eta).
 \label{cp:forward-error}
\end{equation}
The first-layer preactivation error is $a_h$; at layer $j$ it is at
most $9F_{j-1}^h(a_h+\eta)+H_ra_h+A_f\eta$, which proves the displayed
recursion rather than presuming an ambient Lipschitz bound.

Set $d_R=V_R^*w_R-(y-c_n)/\sqrt m$ and $\Delta V=V_C-V_R$.
Subtracting the exact readout formula yields
\begin{equation}
 \widehat w_C-w_R=(I-P_C)z_w-p-T_C\Delta V^*w_R-T_Cd_R.
 \label{cp:readout-deficit}
\end{equation}
Since $p\in\operatorname{ran}V_C$, $(I-P_C)z_w=(I-P_C)\zeta$.
Equations \eqref{cp:readout-observation}--\eqref{cp:forward-error} and
$\|T_C\|\le2/\sqrt\lambda$ imply
\begin{align}
 \|\widehat w_C-w_R\|&\le b_e+6zF(a_h+\eta)
                                      +32zP_h\eta/\sqrt\lambda,
 \label{cp:readout-error}\\
 |f_C(v)-f_n(v)|&\le H_C\{b_e+6zF(a_h+\eta)
                         +32zP_h\eta/\sqrt\lambda\}
                     +3\alpha F(a_h+\eta)+16zP_h\eta.
 \label{cp:output-error}
\end{align}

Let $K_n$ be the dense training tangent Gram: it is the formula
\eqref{cp:selected-gram} with dense fields and the inner products
$n^{-1}u^\top v$ in each neuron layer. The dense deficit therefore obeys
$\dot c_n=-2K_nc_n/m$. The source-only Gram defect
$D=V_R^*V_R+\mathcal J_R^*\mathcal J_R-K_n/m$ satisfies
$\|D\|\le D_s\eta$. To see this without an extra factor $m$, split
each product of two pairings, use \eqref{cp:pair-defect}, and recall
that the sample matrix entries carry the factor $1/m$; a matrix all
of whose entries have magnitude at most $a/m$ has operator norm at
most $a$. With $\Delta\mathcal J=\mathcal J_C-\mathcal J_R$, write
\[
 \Delta\mathcal K=
 V_C^*\Delta V+\Delta V^*V_R+
 \mathcal J_C^*\Delta\mathcal J+
 \Delta\mathcal J^*\mathcal J_R+D.
\]
The error equations are exactly
\begin{align*}
 \dot e&=-2(Q_C+\mathcal J_C^*\mathcal J_C)e
                           -2\Delta\mathcal Kc_n/\sqrt m,\\
 \dot z_w&=2V_Ce+2\Delta Vc_n/\sqrt m,\\
 \dot\zeta&=2\Delta Vc_n/\sqrt m+\dot T_Ce
        -2T_C\mathcal J_C^*\mathcal J_Ce
        -2T_C\Delta\mathcal Kc_n/\sqrt m.
\end{align*}
The term $2V_Ce$ cancels because $T_CQ_C=V_C$.
This is why estimating $c_C-c_n$ by an unconstrained Gronwall
inequality would lose the result.

Here are the needed norm estimates, including the dissipative one.
The derivative formula for $T_C$, and
$\|\dot V_C\|\le10\alpha F_c\rho_C$, give
$\|\dot T_Ce\|\le40zF_c\rho_C\|p\|$.
In the derivative of $\|p\|^2/2$, the $Q_C$ term is $-2u^2$.
The possible absolute contribution from the hidden Gram is at most
$8\|\mathcal J_C\|^2u^2/\lambda\le200z^2F_cu^2$;
the label cap bounds this by $25u^2/(32H_c^2)$, leaving at least
$39u^2/32$. Define
\[
 \mathfrak f=2\|T_C\Delta\mathcal Kc_n/\sqrt m\|,\qquad
 I(t)=\int_0^t\{40zF_c\rho_C b_e+\mathfrak f\}.
\]
Dividing the dissipative inequality by $\|p\|$ gives
\[
 \|p(t)\|+\int_0^t u^2/\|p\|\le I(t),\qquad
 \int_0^t u\le2I(t)/\sqrt\lambda.
\]
At zeros use $(\|p\|^2+\epsilon^2)^{1/2}$ and let
$\epsilon\downarrow0$; $p=0$ implies $e=0$ because $Q_C$ is invertible.
The second inequality also follows directly from
$\|p\|\le2u/\sqrt\lambda$.
Integrating the $\zeta$ equation bounds its hidden-Gram term by
$200z^2F_c I\le25I/32$. Therefore
\begin{equation}
 b_e\le2F\int_0^t\rho_n(a_h+\eta)+3I(t).
 \label{cp:cancelled-energy}
\end{equation}

For completeness, the remaining response estimate can be checked
layer by layer. The dense training preactivation carriers are bounded
coordinatewise by
$M=1+32zK_{\rm src}\sqrt\ell$ on $[0,T]$.
At each gate split
$\phi'(z_C)k_C-\phi'(z_R)k_R$
as $\phi'(z_C)(k_C-k_R)+(\phi'(z_C)-\phi'(z_R))k_R$.
The first multiplication costs $2s$, and the second at most
$2t_2MF_j^z(a_h+\eta)$. Subtracting the initialized and learned reverse
actions costs $A_b\eta$, and the hidden parameter difference costs
$d_{j+1}^E\alpha a_h$. Starting with \eqref{cp:readout-error}, the
displayed $B_j$ recursion, followed by the outer-product split, gives
\[
 \|\Delta\mathcal J\|
 \le B_h\{b_e+M(a_h+\eta)+\eta/\sqrt\lambda\}.
\]
Only the carrier factor contains $M$; it is not multiplied once per
layer. Applying this to the four terms of $\Delta\mathcal K$ gives
\begin{align*}
 \mathfrak f\le{}&2F\rho_n(a_h+\eta)
       +4F\nu(a_h+\eta)/\sqrt\lambda+4D_s\eta\rho_n/\sqrt\lambda\\
 &+4z(J_C+J_R)B_h\rho_n
                  \{b_e+M(a_h+\eta)+\eta/\sqrt\lambda\}.
\end{align*}
The hidden-parameter equation, $\int u\le2I/\sqrt\lambda$ and
$4zJ_C\le5/4$ similarly yield
\[
 a_h\le\tfrac54 I+2B_h\int_0^t\rho_n
                    \{b_e+M(a_h+\eta)+\eta/\sqrt\lambda\}.
\]
Combining with \eqref{cp:cancelled-energy} proves the scalar inequality
\[
 E(t)\le\int_0^t
 \{200zF_c\rho_C+K_1M\rho_n+20F\nu/\sqrt\lambda\}
                \{E+\eta(1+\lambda^{-1/2})\}.
\]
The integral of its coefficient is at most $\mathcal B_n$, by
$\int\rho_n,\int\rho_C\le2z$ and \eqref{cp:source-energy}.
The integrating-factor formula gives
$E\le\eta(1+\lambda^{-1/2})(e^{\mathcal B_n}-1)$.
Substitution in \eqref{cp:output-error} proves the finite-horizon
claim with precisely \eqref{cp:selected-coefficients}.
\end{proof}

\paragraph{Canonical error bound and retained inventory.}
The local fitting claim supplies the tails and endpoint. Thus
\[
 \|f_{\rm model}-f_n\|_{\mathcal X}\le
             \mathcal A_n\eta+\mathcal D e^{-\lambda T/4}.
\]

We verify that the explicit ledger has the promised dependence.
Put $X=\beta^L\ge100$. The local envelopes and fitting bounds give
$u_0,H_c\le X^3$, $\tau\le X^4$ and $K_{\rm src}\le X^{21}$.
For example $16z\le16X^{-30}$ makes both additional terms in
$A_f,A_b$ less than one; hence $A_f,A_b\le19$.
The forward recursion then gives $F\le X^6$.
For further detail, $d_j^E\le X^4$, $H_C\le2X^3$ and
$P_h,P_\delta\le X^4,X^5$. In the backward recursion the additive
term is at most $X^7$, and $18s\le\beta^3$. Geometric summation
over at most $L$ layers gives $\max_jB_j\le X^{11}$ and then
$B_h\le X^{15}$. Also $J_C,J_R\le X^9$ and $D_s\le X^5$.
Substitution gives $K_1\le X^{18}$; here the generous powers
absorb $L\le X$, all numerical factors, and $H_C=2H_c$.
The selected fitting recursion already gives $F_c\le X^{15}$.
Consequently
\[
 \mathcal B_n\le
 400X^{-45}+2X^{-12}+40X^{-24}
                      +64X^{-21}\sqrt\ell
 \le1+\sqrt\ell,
\]
and $\mathcal B_n/z\le C(1+\sqrt\ell)$ with $C$ depending only
on $X$. Since $e^u-1\le ue^u$, \eqref{cp:selected-coefficients}
gives
\[
 \mathcal A_n\le C z(1+\lambda^{-1/2})
                   (1+\sqrt\ell)e^{\sqrt\ell},\qquad
 \mathcal D\le C z(1+\lambda^{-1/2}),
\]
after increasing the activation/depth constant to absorb $e$.
These are exactly the canonical factors in
\eqref{cp:smallcap-comparison}. This is a small-cap proof,
not a claim on a larger recurrence-defined label interval.
Finally $T=32\ell/\lambda$ makes the tail $\mathcal D e^{-8\ell}$.
For fixed positive $Y$ it is at most $Y/(2n)$ eventually; the other
half is \eqref{cp:selected-tolerance}.

\paragraph{Retained inventory.}
Retain both moving and initial copies of the first matrix, mixers and
raw readout; the selected indices; each $M_j$, its inverse or factorization,
and $D_j$; the training data and labels; current training forward and
backward arrays; $c_C$, $V_C$, $Q_C$ and its inverse or factorization;
the corrected readout; direction arrays; and one-query forward workspace.
Every selected width is at most $9q$. Each mixer or metric therefore
costs at most $81q^2$, each first matrix at most $9qd$, each layer's
training arrays at most a fixed multiple of $9mq$, and each sample
matrix at most $m^2$. Since $m,d\le q$, these arrays total at most
$1020(L+1)q^2+36q+10m(d+1)$, including simultaneous initial copies
and working arrays. A nonprimitive activation evaluator is charged
separately. Dense mixers, source bases, coefficient vectors and
coefficient-compilation scratch are not retained by the runtime.
The construction itself uses the finite source spaces only through
its initial arrays; the separate compiler below will establish their
initialization-only provenance for the two methods.
\end{proof}

\subsection{Finite initialized jets suffice}

\begin{proposition}[Initialization-only coefficient compiler]
\label{cp:jets}
On the initialized operator event of \Cref{cp:fit}, fix $T,r_0,M>0$.
Let finitely many sources from \Cref{cp:source} extend coordinatewise
holomorphically near $[-r_0,T+r_0]+i[-r_0,r_0]$, bounded by $M$ uniformly
over their compact query domain. For sphere queries assume also the joint
holomorphic extension of \Cref{cp:source}.
Any finite list of their integrals against specified polynomial time--sphere
bases, or derivatives at real anchors in $[0,T]$, can be approximated to
arbitrary positive accuracy using finitely many derivatives of the dense ODE
and passive recursions at time zero, with exact initialized-image pairing.
No later dense state or trajectory sample is needed; no bound on
initialization time or working storage is asserted.
\end{proposition}
\begin{proof}
We first prove the continuation assertion that makes initialized jets
sufficient. For any one scalar source $g$, use its common bound $M$
on the closed rectangle
$\{-r_0\le\Re t\le T+r_0,\ |\Im t|\le r_0\}$, with a slightly
larger open domain. The source rectangle, or a finite closed-disk cover
with the real interval strictly inside its union, supplies such a
positive $r_0$ by compactness. Put
\[
 \chi=\frac{\pi T}{8r_0},\quad b_0=\tanh\chi,\quad
 b_1=\tanh(\chi+\pi/4),\quad \theta=b_0/b_1,
 \quad \psi(\xi)=\frac{\xi-\theta}{1-\theta\xi},
\]
\[
 \mathfrak t(\xi)=\frac T2+\frac{2r_0}{\pi}
       \log\frac{1+b_1\psi(\xi)}{1-b_1\psi(\xi)},
 \qquad \xi_*=\frac{2\theta}{1+\theta^2}<1.
\]
The logarithm is the branch equal to a real number on the real
interval. For $|\xi|<1$, $|\psi(\xi)|<1$; the ratio in the logarithm
has positive real part and modulus between
$(1-b_1)/(1+b_1)$ and its reciprocal. Thus
$|\Im\mathfrak t|<r_0$ and
$-r_0<\Re\mathfrak t<T+r_0$.
Direct substitution gives $\mathfrak t(0)=0$ and
$\mathfrak t(\xi_*)=T$, and the derivative is nonzero on
$[0,\xi_*]$. This proves continuation from the origin onto the
entire physical interval, without stepping a dense trajectory.

The Taylor coefficients of $g\circ\mathfrak t$ are
\[
 a_j=\sum_{k=0}^j\frac{g^{(k)}(0)}{k!}
                       [\xi^j]\mathfrak t(\xi)^k.
\]
Cauchy's coefficient formula gives $|a_j|\le M$, so truncation at
$N$ has error at most $M\xi_*^{N+1}/(1-\xi_*)$ on the whole interval.
This bound is uniform over all passive inputs in the compact domain
and over all source coordinates. Derivatives at finitely many later
anchors are also recovered: choose $\xi_*<\rho_*<1$ and
$0<\epsilon_*<\rho_*-\xi_*$. On $|\xi|\le\rho_*$ the tail is
$M\rho_*^{N+1}/(1-\rho_*)$, hence its $k$th derivative at any real
anchor is at most
$k!M\rho_*^{N+1}/\{\epsilon_*^k(1-\rho_*)\}$.
Repeated application of $(\mathfrak t')^{-1}\partial_\xi$ converts
these to the finitely many required time derivatives. All its
coefficients are bounded on that compact interval. A finite $N$
therefore achieves every prescribed finite list of accuracies.

Let $\Theta=(W^{(1)},\ldots,W^{(L)},w)$ be the dense parameter tuple
and $\mathcal F$ its autonomous vector field. For the initialized ODE
$\dot\Theta=\mathcal F(\Theta)$ the formal
coefficients are computed successively by
\[
 \Theta_{j+1}=\frac1{j+1}[t^j]
          \mathcal F\!\left(\sum_{i=0}^j\Theta_i t^i\right).
\]
The analytic local solution has precisely these coefficients, by
coefficient comparison and uniqueness of the local ODE. The forward
and passive backward recursions then give their source jets by finite
addition, multiplication and composition at time zero. Any scalar
activation derivatives required here are part of the fixed analytic
activation evaluator; their description, evaluation and scratch are
not asserted to have unit cost unless taken as primitives.

For time--sphere integral coefficients, insert the uniformly convergent
continued polynomials into their defining integrals. Each integrand is
continuous on the compact time interval times sphere. Finite partitions
and weighted Riemann sums converge uniformly for the finite collection
of integrands. A finite mesh can be chosen from their modulus of
continuity: the source strip gives derivative bounds by Cauchy's
formula, and the finitely many basis functions are known polynomials.
Thus the required integral approximations use only finitely many
evaluations of the continued origin-jet polynomial. For $d=1$ the
spatial integral is simply the two-point average.

Finally perform every scalar linear coefficient operation on a base
source and set its image coefficient equal to the initialized matrix
times that vector. Linear operations commute with this fixed matrix.
Since its Euclidean operator norm is at most eight, approximating a
base coefficient in Euclidean norm to $\epsilon/8$ approximates both
the base and its image to $\epsilon$. There are only finitely many
coefficients and a finite sum of their basis suprema; choose each
coefficient tolerance smaller than the desired final error divided
by this sum. This imposes exact pairing while retaining the chosen
uniform source tolerance. No coefficient is retained after the
selected initial arrays are formed.
\end{proof}

\subsection{Whole-sphere polynomial sources}

\begin{lemma}[Joint time--sphere approximation and its dimension]
\label{cp:harmonic-approximation}
Let $T,r_t,r_q,M_n>0$ and $0<\eta\le\min\{1,M_n\}$, and put
$\alpha_T=r_t/(4T)$. Suppose a real $\R^n$-valued source on
$[0,T]\times S^{d-1}$ extends coordinatewise holomorphically to a neighborhood of
\[
 ([-r_t,T+r_t]+i[-r_t,r_t])\times
 \{v\in\mathbb C^d:v^\top v=1,\ \|\operatorname{Im}v\|_2\le\sinh r_q\},
\]
with coordinate modulus at most $M_n$ on this set. Query holomorphy is
relative to the quadric; for $d=1$ use its two real points.
For $d\ge2$ put
\[
 b_d=2d-2,\quad
 D_d=2^{d+1}d^{d-2},\quad
 P_T=\frac{18D_db_d!\,2^{b_d+1}}{\alpha_T r_q^{b_d+1}},
 \quad H_*=2\log\frac{16M_nP_T}{\eta}.
\]
For $\alpha_T,r_q\le1$, a real expansion in Chebyshev time polynomials and
real spherical harmonics with $\alpha_Tk+r_qj\le H_*$ approximates the source
to coordinate error $\eta/8$ using at most
\begin{equation}
 N\le\frac{2\{H_*+\alpha_T+(d-1)r_q\}^d}
                {d!\,\alpha_T r_q^{d-1}}
 \label{cp:harmonic-dimension}
\end{equation}
coefficient vectors. For the dense sources of \Cref{cp:source},
\Cref{cp:jets} gives initialized-jet coefficients with total coordinate
error at most $\eta$ and exact initialized-image pairing.
For $d=1$, use the two points of $S^0$ and at most $2N_1$ vectors,
where $N_1=1+\lfloor H_1/\alpha_T\rfloor$ and
$H_1=\log(64M_n/(\alpha_T\eta))$.
\end{lemma}
\begin{proof}
Fix one coordinate and call its scalar source $g$.
We supply the spatial decay argument. For $d\ge2$ let
$\mathcal H_j$ be the restrictions of homogeneous harmonic
polynomials of degree $j$ to $S^{d-1}$, with probability surface
measure. Its dimension is
\[
 h_j=\binom{j+d-1}{d-1}-\binom{j+d-3}{d-1}
       \le2\binom{j+d-2}{d-2},
\]
where a binomial coefficient with negative upper index is zero.
Indeed, applying the Laplacian to $|x|^{2k}H_l(x)$ gives
$2k(2l+2k+d-2)|x|^{2k-2}H_l(x)$; induction decomposes every
homogeneous polynomial into these summands and gives the dimension
formula by subtraction. The spherical Laplacian eigenvalues
$-j(j+d-2)$ give orthogonality of different degrees by integration
by parts. Polynomials contain constants and separate points on the
compact sphere, so their uniform density (the real Stone--Weierstrass
theorem) proves completeness. For an orthonormal basis $Y_{j\nu}$,
rotational invariance and integration give the addition identity
$\sum_\nu Y_{j\nu}(v)^2=h_j$; consequently
$|Y_{j\nu}(v)|\le\sqrt{h_j}$ and the projection $\Pi_j$ has
supremum operator norm at most $h_j$.

For fixed $v$, average an analytic function over the complex tangent
circle $\cos\zeta\,v+\sin\zeta\,u$, $u\perp v$, $|u|=1$.
On $\mathcal H_j$ this averaging is scalar: rotations fixing $v$
leave a one-dimensional zonal eigenspace, and its normalized
eigenfunction solves
$(1-x^2)p''-(d-1)xp'+j(j+d-2)p=0$, $p(1)=1$.
Coefficient comparison gives
$p(x)=C_j^{(d-2)/2}(x)/C_j^{(d-2)/2}(1)$ for $d>2$.
For real $\zeta$ the averaging is self-adjoint because the joint
law of its two endpoints is symmetric. Apply this to each harmonic
coefficient and continue analytically in $\zeta$; it follows that
the averaged analytic source has $j$th projection equal to
$p(\cos\zeta)\Pi_jg$. At $\zeta=ir_q$ its arguments satisfy
$z^\top z=1$, $|\Im z|\le\sinh r_q$, so its supremum is at most
$M_n$. Thus $\|\Pi_jg\|_\infty\le h_jM_n/p(\cosh r_q)$.

For $\nu=(d-2)/2>0$, the generating function
$(1-2xw+w^2)^{-\nu}=\sum_j C_j^\nu(x)w^j$
at $x=\cosh r$, for any $r>0$, factors as
$(1-e^rw)^{-\nu}(1-e^{-r}w)^{-\nu}$ with nonnegative coefficients.
It gives $C_j^\nu(\cosh r)\ge e^{rj}(\nu)_j/j!$ and
$C_j^\nu(1)=(2\nu)_j/j!$. For $j\ge1$ their ratio is
$2\prod_{i=1}^{j-1}(1+\nu/(\nu+i))$.
Its logarithm is at most
$\log2+\nu\int_0^{j-1}(\nu+x)^{-1}\,dx$.
Since $\nu\ge1/2$, the ratio is at most $2(2j)^\nu$, and hence
at most $2^d(j+1)^d$; the $j=0$ inequality is immediate.
Together with
$h_j\le2d^{d-2}(j+1)^{d-2}$ this proves
\begin{equation}
 \|\Pi_jg\|_\infty\le M_nD_d(j+1)^{b_d}e^{-r_qj}.
 \label{cp:spatial-decay}
\end{equation}
For $d=2$, the basis is the constant and the sine--cosine pair at
each positive integer, and the averaging multiplier is
$\cosh(jr)\ge e^{jr}/2$; the same loose bound with $D_2=8$, $b_2=2$
holds directly.

Set $t=T(1+\cos u)/2$. For $|\Im u|\le\alpha_T\le1$ this lies
inside the time rectangle, since its displacement from $[0,T]$ is
at most $T\alpha_T$. The resulting even periodic function has
cosine coefficients whose $j$th harmonic block is bounded by
$2M_nD_d(j+1)^{b_d}e^{-\alpha_Tk-r_qj}$: shift the Fourier contour
to $\Im u=\pm\alpha_T$ and use \eqref{cp:spatial-decay} there.
The series is absolutely uniformly convergent. Discarding indices
with $\alpha_Tk+r_qj>H_*$ bounds the tail by
\[
 2M_nD_de^{-H_*/2}
 \sum_{k\ge0}e^{-\alpha_Tk/2}
 \sum_{j\ge0}(j+1)^{b_d}e^{-r_qj/2}
 \le M_nP_Te^{-H_*/2}\le\eta/16.
\]
Here $\sum_ke^{-\alpha_Tk/2}\le3/\alpha_T$ and
$\sum_j(j+1)^be^{-r_qj/2}\le3b!(2/r_q)^{b+1}$, obtained by
comparison with the gamma integral after shifting each unit interval.

The number of retained coefficients is
$N=\sum_{r_qj\le H_*}h_j(1+\lfloor(H_*-r_qj)/\alpha_T\rfloor)$.
Replace $h_j$ by twice the number of weak compositions of $j$ into
$d-1$ parts. The corresponding nonnegative lattice points have
weighted sum $\alpha_Tk+r_q\sum l_i\le H_*$. Their disjoint unit
cubes lie in the simplex with right-hand side
$H_*+\alpha_T+(d-1)r_q$, whose volume is exactly the right-hand
side of \eqref{cp:harmonic-dimension}. This proves the count.
For the finite collection of coefficients, choose the Euclidean
approximation tolerance in \Cref{cp:jets} so that its product with
$8\sum_{k,j,\nu}\|Y_{j\nu}\|_\infty$ is at most $\eta/16$.
Both a base source and its paired image then have total error below
$\eta$, without increasing the coefficient count.

If $d=1$, there is no spatial approximation: apply the preceding
one-variable Fourier contour argument separately at the two points.
The geometric tail after degree $\lfloor H_1/\alpha_T\rfloor$ is
at most $6M_ne^{-H_1}/\alpha_T<\eta/8$; the same finite coefficient
accuracy argument finishes this case.
\end{proof}

\begin{proposition}[Harmonic: whole-sphere accuracy and total state]
\label{cp:harmonic}
Eventually at each fixed confidence, the initialized-jet construction gives
a selected nonlinear model with $\|f_{\rm Harm}-f_n\|_{\mathcal X}\le Y/n$
on the whole sphere and total retained real-coordinate count at most
\[
 C(m+d)^2+\left(\frac Cd\right)^{d+1}
       \left(\frac{Ym}\gamma\right)^4[\log(en)]^{3d+2}.
\]
\end{proposition}
\begin{proof}
For this whole-sphere construction put $\lambda=\gamma/m$,
$z=Y/\lambda$, $\ell=\log(en)$, and $T=32\ell/\lambda$.
Choose $\eta$ by the canonical tolerance formula
\eqref{cp:selected-tolerance}. The public source proposition supplies
the radii and coordinate bound
\[
 r_t=\frac{\gamma}
 {\beta^{30L}Y^2m\sqrt{(d+3)\ell}},\qquad
 r_q=\frac1{\beta^{3L}\sqrt{(d+3)\ell}},\qquad
 M_n=\beta^{6L}\sqrt n.
\]
Set $\alpha_T=r_t/(4T)$. For $d\ge2$, let $N$ be the coefficient
count in \Cref{cp:harmonic-approximation}, with its explicitly defined
$P_T,H_*$ evaluated at these local parameters. For $d=1$, let $N_1$
be the two-point time count defined in that same statement.
Apply the approximation lemma to the four joint time--sphere families
from \Cref{cp:source}, including its passive backward family; hence
the same coefficient vectors serve all training inputs. In layer $j$
include the approximate feature and response coefficients and the
paired initialized forward and reverse image coefficients. Add
$\mathbf1_n$, the at most $2m$ exact initialized training vectors
needed for forward-image pairing, and the $d$ first-layer columns.
Each source space has rank at most
\[
 R=2m+d+1+4N\quad(d\ge2),\qquad
 R=2m+d+1+8N_1\quad(d=1).
\]
The source contract is now satisfied. Invoke \Cref{cp:selected} with
its source-dimension bound $q=R$; it selects at most $9R$ coordinates
per layer and gives
the claimed absolute error. All choices use initialized jets only.

We spell out the asymptotic count because dropping the label factor
here would change the theorem. The exact source radius gives
\begin{equation}
 \alpha_T^{-1}=\frac{4T}{r_t}
       =128\beta^{30L}z^2\sqrt{d+3}\,\ell^{3/2},\qquad
 r_q^{-1}=\beta^{3L}\sqrt{(d+3)\ell}.
 \label{cp:time-radius-count}
\end{equation}
For each fixed admissible problem, \eqref{cp:smallcap-comparison}
and \eqref{cp:selected-tolerance} imply
$\log(1/\eta)\le2\ell$ eventually. Every other fixed parameter
inside the logarithm defining $H_*$ is additive, whereas
$\log P_T=O_{d,\mathrm{problem}}(\log\ell)$.
Thus $H_*+\alpha_T+(d-1)r_q\le8\ell$ eventually. The dimension
bound gives
\[
 N\le C^d\beta^{CLd}\frac{(d+3)^{d/2}}{d!}\,
                    z^2\ell^{3d/2+1}.
\]
For $d=1$, the direct two-point count gives the same form with
$N_1\le C\beta^{CL}z^2\ell^{5/2}$.

The public inventory \eqref{cp:selected-inventory}, with $q=R$,
is $1020(L+1)R^2+36R+10m(d+1)$, plus the separately charged
activation evaluator. It already includes selected mixers, metrics,
initial copies, data and working arrays. No dense mixer, source basis,
initialized jet, quadrature table or coefficient vector is retained.

Finally, for $d\ge2$, $R^2\le C(m+d)^2+CN^2$ and
$d!\ge(d/e)^d$ imply
\[
 \frac{(d+3)^d}{(d!)^2}
 \le\frac{(8e^2)^d}{d^{d+1}}.
\]
Indeed, for $d\ge2$, $(1+3/d)^d\le(5/2)^d$ and $d\le2^d$
give the stronger numerator $(5e^2)^d$. Absorb activation/depth
powers into $C^d$ to obtain the asserted count. For $d=1$, use
$R^2\le C(m+1)^2+CN_1^2$ and the already proved bound on $N_1$;
this gives the stated logarithmic exponent five. If a finite small
width makes $9R\ge n$, retaining all coordinates is a valid fallback;
the claimed polylogarithmic bound concerns the eventual range only.
\end{proof}

\subsection{Finite-panel sources on an adaptive time partition}

\begin{proposition}[Taylor: panel accuracy and total state]
\label{cp:panel}
On the declared $(m+p)$-input panel of \Cref{cp:setup}, eventually at each
fixed confidence there is an initialization-only selected model with
$\|f_{\rm Taylor}-f_n\|_{\mathcal X}\le Y/n$ and total state at most
\[
 C(m+p)^2\left[\left(\frac{Ym}\gamma\right)^4[\log(en)]^3+
       \left(\frac m\gamma\right)^2[\log(en)]^2
       \{\log(e+\log(en))\}^2\right]+C(m+p)d.
\]
Only training inputs enter the deficit equation and Gram inverse.
Accuracy outside the panel is not asserted.
\end{proposition}
\begin{proof}
For this finite-panel construction put $\lambda=\gamma/m$,
$z=Y/\lambda$, $\ell=\log(en)$ and $T=32\ell/\lambda$.
Use the public disk-and-partition interface of \Cref{cp:source}:
there are $0=t_0<\cdots<t_J=T$ and radii $r_b>0$ such that each
source is holomorphic near $|t-t_b|\le r_b$ and
$t_{b+1}-t_b\le r_b/2$. Its count \eqref{cp:panel-source} gives
\begin{equation}
 J\le1+\beta^{30L}z^2\sqrt\ell
             +\frac{\beta^{16L}}\lambda
                        \log(1+4\ell/\log2).
 \label{cp:panel-count}
\end{equation}
The centers and radii are deterministic functions of the problem
parameters, as that public interface guarantees, so they require
no later dense training data.

There are at most $2(2m+p)$ source curves per layer: forward features
and their initialized images at all declared inputs, and backward
responses and their initialized transpose images at training inputs.
Each coordinate is bounded by $M_n=\beta^{6L}\sqrt n$ on its disks.
Choose the source tolerance \eqref{cp:selected-tolerance} and
\[
 K=\max\{0,\lceil\log_2(8M_n/\eta)\rceil\}.
\]
On panel $b$, retain coefficient vectors
$r_b^ku^{(k)}(t_b)/k!$ for $0\le k\le K$, where $u$ denotes
one of the source curves.
They have coordinate magnitude at most $M_n$. Since
$|t-t_b|/r_b\le1/2$, the discarded Taylor tail is at
most $M_n2^{-K}\le\eta/8$. The vectors from all panels span one
time-independent space of dimension at most $J(K+1)$ per curve.
The local approximants may differ at a panel boundary: both
approximate the same source there, and the comparison theorem
uses values, never their derivatives or a boundary matching rule.

Apply \Cref{cp:jets} to approximate every scaled base coefficient
in Euclidean norm to $\eta/128$, and form its initialized image
exactly. The base reconstruction error is at most
$(\eta/128)\sum_{k\ge0}2^{-k}=\eta/64$; initialized operator norm
eight gives image error at most $\eta/8$. Adding their independently
controlled Taylor tails leaves both below $\eta$.
Include the constant, exact initial training vectors and paired
images, and first-weight columns required by \Cref{cp:selected}.
Before input
reduction the rank per layer is at most
$2m+d+1+2(2m+p)J(K+1)$.
This verifies the identical source contract, so the same selected
runtime, fitting, cancellation, and endpoint proofs apply. In
particular taking a maximum over the passive inputs does not
introduce $p$ into any comparison coefficient. No passive deficit
or additional sample Gram is introduced.

There is an exact reduction of the first-layer dimension. Let
$U_0\in\R^{d\times k}$ have orthonormal columns spanning all
declared normalized inputs, where $k\le\min(d,m+p)$.
Replace $v_i$ by $U_0^\top v_i$ and the first matrix by
$W^{(1)}U_0$. Every first-layer derivative is a sum of outer
products with $v_a$, hence the component on the orthogonal
complement of this input span is constant. Induction through
the layers proves equality of the original and reduced coupled
trajectories on the panel. The reduced inputs remain unit vectors;
there is no new normalization factor. Each initialized Gaussian
row multiplied by $U_0$ is exactly standard Gaussian in $\R^k$,
and the sample inner products and population feature Gram are
unchanged. Thus the same source event applies to the reduced
system, whose public panel-count coefficients depend only on activations
and depth, not on $k$ or $d$. The source rank is now
\[
 R=2m+k+1+2(2m+p)J(K+1).
\]

Invoke \Cref{cp:selected} with source-dimension bound $q=R$ and
input dimension $k$. To its public inventory
\eqref{cp:selected-inventory}, add the passive inputs and output
buffer and retain the input map. A safe explicit count is
\[
 1020(L+1)R^2+36R+10m(k+1)+pk+(m+p)
                         +3(m+p)d,
\]
plus the separately charged nonprimitive activation evaluator.
This includes original and transformed data, the map, all metrics
and inverse caches, initial and moving weights, direction arrays,
training feature/response and sample-matrix buffers, and query
workspace. Discard the dense arrays, source coefficients, origin
jets, panel anchors and partition schedule. The optimizer is
autonomous; a partition used to prove source approximation is not
a retained forcing table.

For each fixed problem, \eqref{cp:selected-tolerance} implies
$K+1\le8\ell$ eventually. The coefficients $\beta^{30L}$ and
$\beta^{16L}$ in \eqref{cp:panel-count} depend only on activations
and depth. Squaring that bound with
$(a+b+c)^2\le3(a^2+b^2+c^2)$, and using $k\le m+p$, yields
\[
 R^2\le C(m+p)^2\left[
    z^4\ell^3+\lambda^{-2}\ell^2\{\log(e+\ell)\}^2\right].
\]
The second term absorbs the initial rank and plain $\ell^2$ terms
because $\lambda\le1$. Crucially, the $z^2$ in the first term of
the panel count was retained before squaring, so its storage
contribution remains $z^4=(Ym/\gamma)^4$; it was not replaced
by a label-independent constant. This proves the stated inventory
and finishes the construction.
\end{proof}
